%% Commutative diagrams; the arrow heads have to match the font.
\usepackage{tikz-cd}
\tikzcdset{arrow style=math font}

\usetikzlibrary{calc}
\def\centerarc[#1](#2)(#3:#4:#5)% Syntax: [draw options] (center) (initial angle:final angle:radius)
    { \draw[#1] ($(#2)+({#5*cos(#3)},{#5*sin(#3)})$) arc (#3:#4:#5); }
    
%% Shorthands for \mathbb, \mathbf, and \mathcal.  Depending on the TeX version,
%% \bf may still be defined, so we treat this case with a bit more care.
\newcommand{\bb}[1]{\mathbb{#1}}
\newcommand{\ca}[1]{\mathcal{#1}}
\providecommand{\bf}{}
\renewcommand{\bf}[1]{\mathbf{#1}}

%% Typical number sets
\newcommand{\N}{\bb N}
\newcommand{\Z}{\bb Z}
\newcommand{\Q}{\bb Q}
\newcommand{\R}{\bb R}
\newcommand{\C}{\bb C}

%% The opposite of \colon, e.g. in f:X<->Y:g
\newcommand*\lon{%
  \nobreak\mskip4mu plus1mu
  \mathpunct{}%
  \nonscript\mkern-\thinmuskip{:}%
  \mskip1mu\relax}

%% A more efficient way to write spans (by generators)
\newcommand{\gen}[1]{\langle #1\rangle}

%% Some choices for symbols
\AtBeginDocument{
  \let\epsilon = \varepsilon
  \let\le      = \leqslant
  \let\ge      = \geqslant
  \let\leq     = \leqslant
  \let\geq     = \geqslant
}

%% Fixed operators
\newcommand{\Aut}{\mathrm{Aut}}
